\documentclass[11pt]{article}

\usepackage{amsmath, amssymb, amsthm}
\usepackage{geometry}
\usepackage{booktabs}
\usepackage{hyperref}
\usepackage{xcolor}

\geometry{margin=1in}

\newtheorem{proposition}{Proposition}
\newcommand{\code}[1]{\texttt{#1}}

\title{Forward-Looking Expectations and the ICEV Driving Ban:\\
Mathematical Documentation of Model Extensions}
\author{}
\date{}

\begin{document}
\maketitle

\begin{abstract}
\noindent
The base model (as described in the manuscript ``Driving in the wrong
direction? Modelling policy mixes for EV adoption'') gives agents and firms
\emph{naive, permanent-policy} expectations: whatever price or policy applies
today is assumed, in every lifecycle-cost calculation, to hold forever. This
note documents two extensions built on top of that baseline. First, a
\emph{forward-looking expectations} mechanism that replaces the naive
constant-price assumption with a discounted sum over the model's actual,
fully-known future price/policy path, applied consistently to both consumers
and firms, and used in the \code{package/variable\_carbon\_price} experiment
to compare naive and forward-looking responses to a rising carbon price.
Second, an \emph{ICEV driving ban}, implemented by feeding a large cost
shock into that same discounting mechanism rather than as a hand-tuned firm
rule --- so that a forward-looking firm's or consumer's anticipation of the
ban emerges directly from the existing utility function, with no separate
``anticipation lead'' parameter --- used in the \code{package/car\_ban}
experiment to compare ban years 2030--2035. Both extensions are implemented
so that, with the new flags left at their default (off) values, the model's
behaviour is byte-for-byte identical to the pre-existing baseline.
\end{abstract}

\tableofcontents

\section{Baseline: naive, permanent-policy expectations}
\label{sec:baseline}

\subsection{The lifecycle-cost term}

For a vehicle $a$ of transport type ICE or EV, owned or considered by
individual $i$, the model's per-period utility includes a term that
penalises the discounted lifecycle fuel/emissions cost of operating that
vehicle forever. Writing $\omega$ for the vehicle's fuel/energy efficiency,
$\delta$ for its efficiency depreciation rate, $L$ for its current age (in
months), $r$ for the individual's/firm's discount rate, $D_i$ for individual
$i$'s driving distance, $c$ for the per-unit fuel or electricity cost, $e$
for the per-unit emissions intensity, and $\gamma_i$ for individual $i$'s
emissions willingness-to-pay, the existing (pre-extension) formula for this
term is
\begin{equation}
  \mathrm{LCC}_{\mathrm{extra}}
  = D_i \cdot
    \frac{(1+r)(1-\delta)\,\bigl(c + \gamma_i e\bigr)}
         {\omega\,(1-\delta)^{L}\,\bigl(r - \delta - r\delta\bigr)} .
  \label{eq:naive-lcc}
\end{equation}
This is exactly the closed-form expression implemented (identically) in
three places in the code: \code{socialNetworkUsers.py}'s
\code{generate\_utilities\_current}, \code{vectorised\_calculate\_utility\_second\_hand\_cars}
and \code{vectorised\_calculate\_utility\_new\_cars} (the utility a consumer
gets from a vehicle), and \code{firm.py}'s \code{calc\_optimal\_price\_cars} /
\code{calc\_utility} (a firm's own valuation of a vehicle design, used to set
prices and choose what to research/produce).

\subsection{Where Eq.~\eqref{eq:naive-lcc} comes from}
\label{sec:derivation}

Equation~\eqref{eq:naive-lcc} is the closed form of an infinite discounted
sum under two assumptions: (i) the per-unit cost/emissions level
$z \equiv c + \gamma_i e$ is constant forever, and (ii) the vehicle's fuel
efficiency degrades geometrically, $\omega_{L+k} = \omega\,(1-\delta)^{k}$,
$k$ periods after the current age $L$. The present value, as seen at the
vehicle's current age $L$, of all future per-distance operating costs is
\begin{equation}
  C(L) = \sum_{k=0}^{\infty}
    \frac{1}{(1+r)^{k}} \cdot \frac{z}{\omega\,(1-\delta)^{L+k}}
  = \frac{z}{\omega\,(1-\delta)^{L}} \sum_{k=0}^{\infty} q^{k},
  \qquad
  q \equiv \frac{1}{(1+r)(1-\delta)} .
  \label{eq:geom-series}
\end{equation}
Provided $r > \delta/(1-\delta)$ (equivalently $q<1$, the convergence
condition already checked in \code{controller.py}), the geometric series
sums to $1/(1-q)$, and a short calculation gives
\begin{equation}
  \frac{1}{1-q}
  = \frac{1}{1 - \dfrac{1}{(1+r)(1-\delta)}}
  = \frac{(1+r)(1-\delta)}{(1+r)(1-\delta) - 1}
  = \frac{(1+r)(1-\delta)}{r - \delta - r\delta} ,
  \label{eq:series-closed-form}
\end{equation}
using $(1+r)(1-\delta) - 1 = r - \delta - r\delta$. Substituting
\eqref{eq:series-closed-form} into \eqref{eq:geom-series} and multiplying by
$D_i$ reproduces Eq.~\eqref{eq:naive-lcc} exactly. This is the calculation
this document refers to as the \emph{naive/permanent} assumption: it
correctly discounts the future, but it assumes the price/emissions level $z$
observed \emph{today} will never change --- so a rising carbon price, for
instance, is invisible to it until the month it actually arrives.

\section{Forward-looking expectations}
\label{sec:forward-looking}

\subsection{Generalising the geometric series}

The paper itself notes this limitation directly (Section~4.1): ``we
disregard the role of expectations regarding policy stringency \ldots\ with
agents assuming that policies will be permanent.'' The fix does not require
abandoning Eq.~\eqref{eq:geom-series}'s structure, only its assumption that
$z$ is constant. Write $z_t$ for the per-unit cost or emissions level at
absolute simulation month $t$ (known for the whole horizon in advance,
since it is simply the pre-computed price/carbon-price/subsidy path already
built by \code{controller.py}). Replacing the constant $z$ in
\eqref{eq:geom-series} by the actual path $z_{t+k}$ gives the
forward-looking present-value index
\begin{equation}
  \mathrm{Index}(t) \;\equiv\; \sum_{k=0}^{T_{\max}-t} z_{t+k}\, q^{k},
  \qquad q = \frac{1}{(1+r)(1-\delta)},
  \label{eq:index-def}
\end{equation}
where $T_{\max}$ is the last simulated month. When $z_t = z$ for all $t$,
\eqref{eq:index-def} collapses to $z/(1-q)$, i.e. exactly the naive case ---
so Eq.~\eqref{eq:index-def} is a strict generalisation of
Eq.~\eqref{eq:geom-series}, not a different model.

\subsection{Backward recursion}

Rather than evaluating the sum in \eqref{eq:index-def} from scratch at every
$t$ (an $O(T_{\max}^2)$ computation, and one that would have to be repeated
per agent if it were not for the separation argument in
Section~\ref{sec:separation}), it is computed once, for the whole horizon,
via the backward recursion
\begin{equation}
  \mathrm{Index}(t) = z_t + q\,\mathrm{Index}(t+1),
  \qquad
  \mathrm{Index}(T_{\max}) = \frac{z_{T_{\max}}}{1-q},
  \label{eq:backward-recursion}
\end{equation}
which is easily verified by unrolling \eqref{eq:index-def} one step:
$\sum_{k=0}^{K} z_{t+k} q^k = z_t + q \sum_{k=0}^{K-1} z_{t+1+k} q^{k}$. The
boundary condition treats the last known value $z_{T_{\max}}$ as persisting
forever beyond the modelled horizon --- the same perpetuity assumption the
naive model already makes everywhere, just deferred to ``after the
simulation stops knowing anything'' rather than applied from today onward.
This is implemented as \code{Controller.compute\_discounted\_indices()} /
\code{Controller.\_discounted\_index()}, a single $O(T_{\max})$ pass computed
once per model run (or once per future-period re-run when reusing a
calibration, see \code{setup\_continued\_run\_future}), independent of the
number of agents or vehicles.

\subsection{Why cost and emissions must be indexed separately}
\label{sec:separation}

The naive formula \eqref{eq:naive-lcc} combines cost and emissions
\emph{before} discounting: $z = c + \gamma_i e$. This is harmless in the
naive case because $\gamma_i$, although agent-specific, is just a constant
multiplier sitting outside an otherwise-agent-independent sum. It stops
being harmless once $c_t$ and $e_t$ are allowed to vary independently over
time (as they do once carbon price, gasoline price and grid-decarbonisation
paths differ): $\gamma_i$ cannot be pulled inside a sum that is shared
across all agents if the combination $c_t + \gamma_i e_t$ is what gets
discounted, since that would require recomputing the whole discounted sum
once per distinct $\gamma_i$ value. Instead, the implementation keeps
\emph{two} indices,
\begin{equation}
  \mathrm{Index}_{\mathrm{cost}}(t) = \sum_{k} c_{t+k}\, q^{k},
  \qquad
  \mathrm{Index}_{\mathrm{emissions}}(t) = \sum_{k} e_{t+k}\, q^{k},
  \label{eq:two-indices}
\end{equation}
each computed once (globally, not per-agent, since $r$ and $\delta$ are
global scalars shared by all ICE, resp.\ all EV, vehicles), and combines
them with the agent-specific $\gamma_i$ \emph{outside} the sum, at the point
of use:
\begin{equation}
  \mathrm{LCC}_{\mathrm{extra}}^{\mathrm{FL}}
  = D_i \cdot
    \frac{\mathrm{Index}_{\mathrm{cost}}(t) + \gamma_i\, \mathrm{Index}_{\mathrm{emissions}}(t)}
         {\omega\,(1-\delta)^{L}} .
  \label{eq:fl-lcc}
\end{equation}
This is the forward-looking replacement for Eq.~\eqref{eq:naive-lcc}. Note
the age-decay factor $\omega\,(1-\delta)^{L}$ is unchanged: it depends only
on the vehicle's efficiency and current age, not on the price/emissions
path, so it factors out of the sum in \eqref{eq:geom-series} identically in
both the naive and forward-looking cases. For newly-designed cars (age
$L=0$, as evaluated by firms in Section~\ref{sec:firms}), this factor is
simply $\omega$.

\subsection{The four price/emissions paths}

Four scalar paths are discounted via \eqref{eq:backward-recursion}, one
pair (cost, emissions) for each of the two energy carriers:
\begin{align}
  z^{\mathrm{gas}}_{\mathrm{cost}}(t)
    &= \mathrm{GasPrice}(t) + \mathrm{CarbonPrice}(t)\cdot \bar e_{\mathrm{gas}},
    &
  z^{\mathrm{gas}}_{\mathrm{em}}(t) &= \bar e_{\mathrm{gas}}
    \quad (\text{constant: ICE tailpipe intensity does not change over time}),
  \\
  z^{\mathrm{elec}}_{\mathrm{cost}}(t)
    &= \mathrm{ElecPrice}(t)\,\bigl(1-s(t)\bigr) + \mathrm{CarbonPrice}(t)\cdot \bar e_{\mathrm{grid}}(t),
    &
  z^{\mathrm{elec}}_{\mathrm{em}}(t) &= \bar e_{\mathrm{grid}}(t),
  \label{eq:four-paths}
\end{align}
where $s(t)$ is the electricity price subsidy proportion and
$\bar e_{\mathrm{grid}}(t)$ is the (time-varying, decarbonising) grid
emissions intensity. These are exactly \code{gas\_price\_california\_vec},
\code{carbon\_price\_time\_series}, \code{gas\_emissions\_intensity},
\code{electricity\_price\_vec}, \code{electricity\_price\_subsidy\_time\_series}
and \code{electricity\_emissions\_intensity\_vec} already built by
\code{controller.py}'s existing calibration/scenario/policy machinery ---
\textbf{no new exogenous path had to be built}; only the mechanism that
reads it changed. Each of the four resulting index arrays
($\mathrm{Index}^{\mathrm{gas}}_{\mathrm{cost}}$,
$\mathrm{Index}^{\mathrm{gas}}_{\mathrm{em}}$,
$\mathrm{Index}^{\mathrm{elec}}_{\mathrm{cost}}$,
$\mathrm{Index}^{\mathrm{elec}}_{\mathrm{em}}$) is looked up at the current
month and written onto every car of the matching transport type as
\code{car.cost\_index} / \code{car.emissions\_index}, exactly mirroring how
\code{car.fuel\_cost\_c} / \code{car.e\_t} are already refreshed every
timestep.

\subsection{Firms}
\label{sec:firms}

The identical substitution (Eq.~\eqref{eq:naive-lcc}
$\to$ Eq.~\eqref{eq:fl-lcc}) is applied on the firm side, in
\code{Firm.\_lifecycle\_cost\_term\_firm}, used both when a firm prices a
car for sale (\code{calc\_optimal\_price\_cars}) and when it values a
candidate design during R\&D (\code{calc\_utility}, called from
\code{calc\_utility\_cars\_segments} and
\code{set\_utility\_and\_profit\_grid\_production}). Since a firm always
prices/evaluates a freshly designed car, its age is $L=0$ and the age-decay
factor is simply $\omega$ (i.e. \code{age\_factor=1} in the code, verified
to be bit-for-bit equivalent to omitting the factor entirely).

\subsection{A single switch, two exact regimes}

Everywhere above, the choice between Eq.~\eqref{eq:naive-lcc} and
Eq.~\eqref{eq:fl-lcc} is governed by one boolean,
\code{forward\_looking\_expectations}, read once per run (default
\code{False}) and shared by every consumer and every firm in that run. When
it is \code{False}, the code path is Eq.~\eqref{eq:naive-lcc}, expressed
with the exact same operations/operand order as the original code --- so
existing configurations that never set this flag get bit-for-bit identical
output. When it is \code{True}, every consumer and firm uses
Eq.~\eqref{eq:fl-lcc}, driven by the same four index arrays.

\subsection{Application: the \code{variable\_carbon\_price} experiment}

\code{package/variable\_carbon\_price} exploits Eq.~\eqref{eq:fl-lcc}
directly: the carbon-price path $\mathrm{CarbonPrice}(t)$ is set to $\sim 10$
different linearly-rising ramps (reaching a range of end-prices at the end
of the policy horizon) and to $5$ flat carbon taxes at correspondingly
\emph{half} the ramps' end-prices (so that a flat tax and its reference ramp
have the same time-averaged carbon price --- a linear ramp from $0$ to $X$
averages $X/2$, so comparing it to a flat tax of $X$ would overstate how
much cheaper the ramp looks). Every scenario is run under both
$\code{forward\_looking\_expectations}\in\{\code{False},\code{True}\}$. Under
naive expectations, a rising carbon price is met reactively, one month at a
time, as though each month's price were permanent; under forward-looking
expectations, Eq.~\eqref{eq:index-def} means agents and firms price in the
\emph{entire future ramp} from the start, so EV adoption should shift
earlier and be less sensitive to the exact shape of the ramp, relative to
the naive case. This is exactly the comparison the experiment's naive-vs-FL
panels are designed to surface.

\section{ICEV driving ban}
\label{sec:ban}

\subsection{A first attempt, and why it was rejected}
\label{sec:ban-first-attempt}

An earlier version of this feature modelled a \emph{sale} ban: firms could
no longer sell new ICE cars from a fixed date, and a forward-looking firm
was given a separate, hand-chosen ``anticipation lead'' parameter that made
it stop \emph{researching} new ICE models some number of months early. This
was rejected for a good reason: the firm's R\&D/production decision in this
model (\code{calc\_predicted\_profit\_segments\_research}/\code{\_production}
in \code{firm.py}) is a \emph{single-period} profit maximisation --- it has
no notion of amortising R\&D cost over a design's future selling life, so
there was no mechanism by which ``fewer years left to sell this car'' could
naturally make it look less profitable. Making the wind-down emerge from
profitability alone would have required building genuinely new multi-period,
discounted-profit machinery for firms. The anticipation lead was, instead,
an arbitrary dial with no basis in the rest of the model --- exactly the
kind of naive/permanent-vs-forward-looking asymmetry this whole extension
is meant to avoid introducing by hand.

\subsection{A driving ban lets anticipation emerge from the existing utility function}

The fix is to change \emph{which} policy is being modelled, not to patch
the firm decision rule. A \textbf{driving ban} --- from $T_{\mathrm{ban}}$
onward, ICE cars become prohibitively costly to operate at all, not merely
unavailable for new sale --- can be expressed as a cost shock, and a cost
shock is exactly what Section~\ref{sec:forward-looking} already knows how
to discount. Concretely, $z^{\mathrm{gas}}_{\mathrm{cost}}(t)$ from
Eq.~\eqref{eq:four-paths} is redefined as
\begin{equation}
  z^{\mathrm{gas}}_{\mathrm{cost}}(t)
  = \mathrm{GasPrice}(t) + \mathrm{CarbonPrice}(t)\cdot \bar e_{\mathrm{gas}}
    \;+\; B \cdot \mathbb{1}[t \ge T_{\mathrm{ban}}],
  \label{eq:ban-cost-shock}
\end{equation}
where $B$ (\code{ICE\_driving\_ban\_penalty}) is a constant added to the
effective per-unit ICE fuel cost for every month from $T_{\mathrm{ban}}$
onward, on top of whatever gas/carbon-price cost already applied, and fed
into the \emph{same} $\mathrm{Index}_{\mathrm{cost}}$ machinery used for the
carbon price. Crucially, this requires \textbf{no new mechanism and no
firm-specific code whatsoever}:
\begin{itemize}
  \item \textbf{Naive agents/firms} read only $z_t$ at the current month
    (Eq.~\eqref{eq:naive-lcc}), so they are unaffected until $t$ actually
    reaches $T_{\mathrm{ban}}$, at which point Eq.~\eqref{eq:naive-lcc}'s
    existing amplification factor $(1+r)(1-\delta)/(r-\delta-r\delta)
    \approx 1/(1-q)$ turns even a moderate $B$ into an overwhelming utility
    penalty, and the ordinary logit choice mechanism does the rest --- ICE
    is abandoned rapidly, but only \emph{after} the ban takes effect.
  \item \textbf{Forward-looking agents/firms} discount the whole known
    future path (Eq.~\eqref{eq:index-def}), so $\mathrm{Index}_{\mathrm{cost}}(t)$
    rises smoothly as $t\to T_{\mathrm{ban}}$ even while $t < T_{\mathrm{ban}}$
    --- consumers shift away from ICE, and firms' expected profit from ICE
    erodes right along with the collapsing demand, so
    \code{choose\_cars\_segments()}/\code{innovate()} shift toward EV on
    their own, entirely through the existing profit-maximisation logic. No
    anticipation-lead parameter is chosen anywhere; how far in advance this
    is felt is governed purely by $q=1/((1+r)(1-\delta))$, i.e. by the
    model's own calibrated discount rate and depreciation rate.
\end{itemize}
Because Eq.~\eqref{eq:ban-cost-shock} feeds the same fuel-cost channel that
\code{secondHandMerchant.py} already reads every step, the ban also reaches
second-hand ICE cars automatically --- correct for a driving ban, which
(unlike a sale ban) makes every ICE car on the road illegal to drive
regardless of when or from whom it was bought.

The emissions path is deliberately left untouched
($z^{\mathrm{gas}}_{\mathrm{em}}(t) = \bar e_{\mathrm{gas}}$ throughout): the
ban penalty represents illegality/unavailability, not a change in the
physical emissions of a car that is, in the small residual probability the
logit choice model still assigns it, still driven.

\subsection{Bans beyond the simulated horizon}
\label{sec:ban-horizon}

A subtlety in a naive implementation of Eq.~\eqref{eq:ban-cost-shock}: the
arrays underlying $\mathrm{Index}_{\mathrm{cost}}$ only span the simulated
horizon, $n = $ \code{duration\_burn\_in + duration\_calibration +
duration\_future} months. If $T_{\mathrm{ban}} \ge n$ (e.g. a ban configured
for 2050 on a run that only simulates to 2030), a version of
\eqref{eq:backward-recursion} restricted to $t<n$ never actually evaluates
$z_t$ for any $t\ge T_{\mathrm{ban}}$, so the ban would have \emph{no effect
whatsoever} — silently identical to BAU — which contradicts the premise
that a forward-looking agent has full knowledge of the future policy path
regardless of how far ahead the simulation happens to run.

The fix extends the \emph{recursion}, not the stored arrays actually read
during simulation, out to $\max(n,\, T_{\mathrm{ban}}+1)$ whenever
$T_{\mathrm{ban}}\ge n$, holding gas and carbon price constant over the
extension — the same perpetuity assumption
\eqref{eq:backward-recursion}'s own boundary condition already makes at
$t=T_{\max}$ — before slicing the result back down to the first $n$
entries. Verified numerically on a 372-month run (duration\_future
truncated to 72 months, i.e. simulating only to $\sim$2030): with a ban
configured 240 months beyond the horizon (a 2050-equivalent date), the
forward-looking cost index at the very start of the future period is
$10{,}561$ against a BAU value of $54$ and a within-horizon
(2030-equivalent) ban's $18{,}364$ at the same reference point ---
confirming the distant ban is felt (dramatically more than BAU), and felt
\emph{less} than a nearer one, exactly the ordering a rational
infinite-information agent should produce.

\subsection{An emergent, and informative, finding}
\label{sec:ban-emergent-finding}

Because $\delta \approx 0.00175$ and $r\approx 0.00407$ in the calibrated
model, $q = 1/((1+r)(1-\delta)) \approx 0.9977$, giving an effective
discounting horizon $1/(1-q)\approx 433$ months ($\sim$36 years). A direct
consequence, verified numerically: for a driving ban 6--11 years out (the
2030--2035 range this experiment sweeps), Eq.~\eqref{eq:index-def}'s
backward recursion means $\mathrm{Index}_{\mathrm{cost}}(t)$ is already
70--85\% of its post-ban plateau value from essentially the \emph{start} of
the simulated future period, for \emph{every} ban year in that range ---
not a small, late-arriving nudge. This is not an artefact of the
implementation; it is what the model's own long effective planning horizon
implies once agents are given perfect information about the future. One
practical upshot is that forward-looking outcomes differ from naive
outcomes almost immediately in this experiment, and differ from each other
across ban \emph{years} comparatively mildly (Section~\ref{sec:ban-application}) ---
a genuine, load-bearing implication of the calibration, not a tuning choice
made for this feature.

\subsection{Application: the \code{car\_ban} experiment}
\label{sec:ban-application}

\code{package/car\_ban} sweeps $T_{\mathrm{ban}}$ over calendar years
2030--2035 (converted to absolute months via the same reference point the
existing \code{controller.t\_2030} constant uses, i.e. the future period
begins at calendar year 2024), each under both
$\code{forward\_looking\_expectations}\in\{\code{False},\code{True}\}$, plus
a no-ban (BAU) baseline. A smoke test at small scale (2 seeds, a single
2030 ban year) confirms the mechanism behaves as designed: with
$B=50$ (against a calibrated gas price of order 1--2 per unit), naive
EV fleet share tracks BAU right up to the ban and only drifts upward
afterward (0.40 at the ban, 0.51 six months later, 0.74 two years later),
while forward-looking EV fleet share under the identical ban reaches 0.97
\emph{six months before} the ban takes effect, and 0.95 a full year before
it --- versus 0.39 for naive agents facing the same ban at the same
horizon. Comparing the naive and forward-looking columns of each metric's
figure shows what anticipation of the ban, rather than the ban itself
(which every regime experiences identically once $t\ge T_{\mathrm{ban}}$),
changes about the transition.

\section{Summary of new parameters}

\begin{table}[h]
\centering
\begin{tabular}{@{}lll@{}}
\toprule
Parameter & Default & Meaning \\
\midrule
\code{forward\_looking\_expectations} & \code{False} & Eq.~\eqref{eq:naive-lcc} vs. Eq.~\eqref{eq:fl-lcc} \\
\code{ICE\_driving\_ban\_time} & \code{None} (never) & $T_{\mathrm{ban}}$, Eq.~\eqref{eq:ban-cost-shock} \\
\code{ICE\_driving\_ban\_penalty} & $0$ & $B$, Eq.~\eqref{eq:ban-cost-shock} \\
\bottomrule
\end{tabular}
\caption{New backward-compatible controller parameters. All three default
to values that reproduce the pre-existing model exactly.}
\end{table}

\end{document}
